---
date: ``2026-07-30''
tags:
  - daily
draft: ``false''
aliases:
---

up:: 011b MOC Capital II

\section{Cap. 8: ``Capital fixo e capital circulante''}
Distinção entre capital fixo e capital circulante surge na análise das rotações do capital. 
\begin{quote}
``Essa cessão de valor ou transferência do valor de tal meio de produção ao produto, para cuja criação ele colabora, é determinada por um cálculo médio: pela duração média de sua função desde o momento em que o meio de produção entra no processo de produção até o momento em que ele está completamente desgastado, morto, e tem de ser substituído ou reproduzido por um novo exemplar do mesmo tipo.'' \parencite[p. 239-40]{Marx2014}
\end{quote}

Dessa forma, tal distinção é apenas no que perdura ou não em forma de valor de uso, pois, em particular no caso do capital fixo, ``é apenas seu valor que circula e, mesmo assim, gradualmente, de modo fragmentado, à medida que vai sendo transferido ao produto, que circula como mercadoria'' (p. 240-1).

Assim como um mesmo objeto pode ter valores de uso distintos, da mesma forma um meio de trabalho pode ter função distinta como capital fixo ou circulante:
\begin{quote}
``O gado, como gado de trabalho, é capital fixo; como gado de corte, é matéria-prima, que entra na circulação como produto, ou seja, como capital circulante, não fixo.'' \parencite[p. 244]{Marx2014}
\end{quote}

Em suma, a distinção entre capital fixo e circulante:
\begin{quote}
``2. A rotação dos elementos fixos que formam o capital e, portanto, também o tempo de rotação aí requerido, abarca várias rotações da parte líquida [circulante] do capital. A cada rotação do capital fixo, o capital líquido [circulante] realiza várias rotações. Uma das partes integrantes do valor do capital produtivo só assume a forma específica do capital fixo quando o meio de produção que constitui esse valor \textit{não} se consome durante o tempo de confecção do produto e sai do processo de produção transformado em mercadoria. Uma parte de seu valor tem de permanecer vinculada à antiga forma de uso, enquanto a outra parte é posta em circulação pelo produto acabado, cuja circulação faz circular, ao mesmo tempo, o valor total dos elementos líquidos que formam o capital.'' \parencite[p. 249, grifo meu]{Marx2014}
\end{quote}

Além disso, o capital fixo requer que seja adiantado ``de uma só vez'', mas só tem seu valor retirado novamente da circulação ``de modo fragmentado e gradual'' (i.e. depreciação) (p. 250). Em contraste, o capital circulante — constante e capital variável — tem sempre de ser renovado ``em espécie'', ``\textit{in natura}'', ``os meios de produção [circulantes] por novos exemplares do mesmo tipo, a força de trabalho por sua compra sempre renovada'' (p. 250). O caso da força de trabalho é diferente, pois ``encontra-se constantemente no processo de produção, mas apenas por meio da \textit{renovação constante de sua compra}, e \textit{frequentemente com a troca das pessoas}.'' (p. 250, grifo meu).

\subsubsection{Detalhes sobre capital fixo}
\begin{quote}
``Na determinação do desgaste e dos custos de reparo segundo uma média social, resultam necessariamente grandes desigualdades, mesmo tratando-se de investimentos de capital de mesmo volume, efetuados sob as mesmas circunstâncias e no mesmo ramo de produção. Na prática, o que ocorre é que, para um capitalista, a máquina etc. dura mais do que seu período médio de vida, para outro, ela dura menos. Os custos de reparo do primeiro ficam acima, e os do segundo ficam abaixo da média etc. Mas o acréscimo ao preço da mercadoria, determinado pelo \textit{desgaste e pelos custos de reparo} desta última, \textit{é o mesmo e é determinado por aquela média social}. Um deles, portanto, obtém com esse acréscimo mais do que ele realmente acrescenta, e o outro menos.'' \parencite[p. 260]{Marx2014}
\end{quote}

A depender do tempo médio de vida de um capital, é possível que, na prática,
\begin{quote}
 ``o desgaste e, portanto, também sua reposição, torn[e]-se uma grandeza praticamente insignificante, de modo que somente os custos de reparo são computados.'' \parencite[p. 263]{Marx2014}
\end{quote}

Imagino que seja o caso, por exemplo, de fiações elétricas, usinas hidrelétricas, etc. (Mercadoria-Energia.).

\begin{quote}
``Isso vale para todas as obras de duração secular, em que, portanto, o capital nelas investido não precisa ser gradualmente reposto na proporção de sua depreciação, bastando que os custos médios anuais de manutenção e reparo sejam transferidos ao preço do produto.'' \parencite[p. 264]{Marx2014}
\end{quote}

(Aqui, Marx está falando ``na prática'', num nível mais concreto. Evidentemente que preços de mercado tenderão a comportar tais custos de depreciação em seus produtos... ou ao menos espera-se.)

\section{Cap. 9: ``A rotação total do capital desembolsado. Ciclos de rotação''}
\begin{quote}
``1. A rotação total do capital desembolsado é a rotação média de seus diversos componentes'' \parencite[p. 267]{Marx2014}
\end{quote}

O tempo de rotação do capital é uma média harmônica ponderada dos tempos de rotação de seus componentes. Tal ``harmonização'' dos tempos distintos de rotação das componentes de um capital, que possuem características qualitativamente distintas entre si, requer sua redução ``a uma forma homogênea de rotação, de modo que elas se diferenciem apenas quantitativamente, segundo a duração de suas rotações'' (p. 268) — evidentemente precisamos analisá-las segundo o Ciclo do Capital-Dinheiro. 

Embora o objetivo do capital seja o de sua autovalorização o mais rápido possível, i.e. de diminuir seu tempo de rotação o máximo que puder, convém-lhe empregar mais e mais capital fixo justamente para que
\begin{quote}
``o valor de capital que efetue sua rotação \textit{durante o ano} pode, em consequência das repetidas rotações do capital líquido [circulante] durante esse mesmo ano, ser maior que o valor total do capital desembolsado [i.e. adiantado].'' (p. 268, grifo meu)
\end{quote}

Ou seja, é justamente porque o emprego da maquinaria permite um aumento sobremaneira das forças produtivas — leia-se: de maior produção material (maior quantidade) em dado tempo — que ele torna-se crucial à autovalorização célere do capital. Isso é justamente porque ``A \textit{rotação do valor} do capital desembolsado se separa (...) de seu período real de reprodução ou do período real de rotação de seus componentes.'' (p. 269).

Além disso, também há, devido à constante revolução técnica-tecnológica dos meios de produção — i.e. de ganhos de produtividade —, o surgimento de depreciação moral, i.e. da obsolescência prematura destes meios em face de maquinaria mais produtiva no mercado, a qual aumenta ``a necessidade de sua constante reposição (...) muito antes que estejam esgotados fisicamente'' (p. 269).

Neste caso, porém, o tempo de rotação é uma mera média harmônica dos tempos de rotação das componentes do capital pois supomos que não há reprodução ampliada, i.e. se há mais-valor produzido, ele não é capitalizado (no que diz respeito a este capital). A produção de mais-valor permite com que se acelere o tempo de rotação do capital, ao retornar o valor adiantado mais rápido.

\section{Capítulo 12: ``O período de trabalho''}
Abstraindo de tempo de circulação e de tempo ocioso ($T_{\lnot trab}$), temos que o tempo de rotação é dado totalmente pelo tempo de trabalho, que será o foco deste capítulo.
\begin{quote}
``Com igual adiantamento de capital, a diferença na duração do ato da produção tem de provocar, evidentemente, uma diferença na velocidade da rotação, ou seja, nos períodos para os quais um dado capital foi adiantado.'' \parencite[p. 320]{Marx2014}
\end{quote}

O que Marx denomina Período de Trabalho é a ``grandeza contínua (...) formada pela sucessão de um número maior ou menor de jornadas de trabalho conexas'' (p. 321).
\begin{quote}
``A uma tal jornada de trabalho, formada pela sucessão de um número maior ou menor de jornadas de trabalho conexas, denomino \textit{período de trabalho}. (...) isto significa o número de jornadas de trabalho conexas que, num ramo determinado de negócios, é requerido para fornecer um produto acabado. O produto de cada jornada de trabalho é aqui apenas um \textit{produto parcial}, que se segue executando diariamente e que só recebe sua figura acabada, \textit{só é um valor acabado de uso ao fim de um período mais ou menos prolongado de tempo de trabalho}.'' \parencite[p. 321, grifo meu]{Marx2014}

``Durante todo o período de trabalho vai-se acumulando, em camadas, a parte de valor que o capital fixo transfere diariamente ao produto até este atingir sua forma acabada.'' (p. 322)
\end{quote}

Aqui se faz visível, na prática, o que significa algo ser capital fixo: 
\begin{quote}
``\textit{Enquanto cada período singular} no interior do qual seu valor reflui fracionariamente mediante a venda do produto \textit{for mais curto que seu próprio período de existência}, a mesma máquina a vapor continuará a funcionar no processo de produção \textit{durante vários períodos de trabalho}.'' \parencite[p. 322, grifo meu]{Marx2014}
\end{quote}

Por exemplo, uma máquina que tem duração média (para dado uso industrial etc.) de $20$ anos, então ela poderá ser empregada por vários períodos de trabalho em e.g. uma indústria têxtil que produz diariamente e a cuja produção ``transfere seu valor fracionada e diariamente'', tanto quanto em uma indústria ferroviária em que transfere valor ``ao longo de três meses''. (ibid.).

A importância do tempo de rotação é que, quanto menor ele for, tanto mais rápido o valor nele despendido refluirá ao capitalista; ou seja, tão mais rápido ``pode-se comprar nova força de trabalho e novas matérias-primas'' (p. 323). O oposto se dá em setores com produção mais ``perene'', com tempo de rotação maior (p. ex. ferrovias etc.), em que o tempo de refluxo do valor é maior, portanto requerendo mais capital adiantado para a manutenção dos períodos de trabalho\footnote{Além de Custos de Circulação, ``\textit{faux frais}'' que não entram no valor da mercadoria. } até o momento da venda do produto e o refluxo do valor total até ali despendido.
\begin{quote}
``num caso o mesmo capital se renova rapidamente e, assim, pode repetir a mesma operação, ao passo que, no outro, ele só se renova de maneira relativamente lenta, de modo que, até o prazo de sua renovação, novas quantidades de capital têm de ser constantemente adicionadas às quantidades antigas.'' (p. 324)
\end{quote}

Aqui já me remete a todo o conceito de ``finanças pacientes'' e de investimentos verdes etc. Me parece algo bem antitético ao capital....... mas então por que é que ainda tem capitalistas que investem neles, se for o caso!? Eles são tão rentáveis assim!? Ou quem sabe são justamente ramos em que capital-dinheiro ocioso de outros ramos possa ser despendido!

Resposta já está na mesma página: tais investimentos de prazo mais longo e que requerem muito capital adiantado, "Nas fases menos desenvolvidas da produção capitalista, 
\begin{quote}
``ou não são em absoluto executados de modo capitalista, como é o caso, por exemplo, de estradas, canais etc., construídos \textit{à custa da comunidade ou do Estado}'' (como é o caso prototípico da definição de Bens Públicos da economia usual), ``(...) ou, então, (...) só são fabricados numa quantidade ínfima e custeados \textit{pelo próprio patrimônio do capitalista}.'' (p. 324, grifo meu)
\end{quote}

Com o desenvolvimento da sociedade capitalista, diminui-se sobremaneira essas quantidades mínimas de capital necessárias para o que outrora eram considerados investimentos de magnitude e escala temporal monumentais, tanto pelas ``somas maciças de capital [que] se concentram nas mãos de indivíduos'' quanto por conta das sociedades por ações e pelo sistema mais desenvolvido de crédito. \[A ``financeirização'' é um reflexo disso, em que o Capital Fictício não só se espraiou sobre todas as atividades do capitalismo — até um bom ponto, já o era na própria época de Marx —, quanto permite a captação de somas maiores ou menores de Dinheiro/mercadoria-capital.\]

Ou seja, tais empreitadas só passam a integrar plenamente a produção capitalista 
\begin{quote}
``quando a concentração de capital já é bastante considerável e, ao mesmo tempo, o desenvolvimento do sistema de crédito oferece ao capitalista o cômodo expediente de adiantar — e, desse modo, também arriscar — capital alheio, em vez de capital próprio'' \parencite[p. 326]{Marx2014}
\end{quote}

\subsubsection{O ônus e o bônus do aumento das forças produtivas via capital fixo}
\begin{quote}
``As circunstâncias que aumentam o produto de cada jornada de trabalho individual, como a cooperação, a divisão do trabalho, o emprego da maquinaria, encurtam ao mesmo tempo o período de trabalho nos atos conexos de produção. (...) No entanto, esses aperfeiçoamentos, que encurtam o período de trabalho e, por conseguinte, o tempo pelo qual o capital circulante tem de ser adiantado, estão geralmente vinculados a um desembolso aumentado de capital fixo.'' 
\parencite[p. 326]{Marx2014}
\end{quote}

\section{Cap. 13: ``O tempo de produção''}
Nem todo tempo de produção é tempo de trabalho. Pode haver, e necessariamente há, um período de tempo no qual o produto inacabado ``permanece fora do processo de trabalho, deixad[o] à ação de processos naturais'' \parencite[p. 332]{Marx2014}.

Este tempo de ``não-trabalho'' pode ser passível de diminuições, por exemplo uma maior intensidade do trabalho, em que uma jornada de trabalho se torna ``menos porosa'', i.e. em que se ``reduz o dispêndio improdutivo de meios de produção, matérias-primas e força de trabalho'', com o que ``se reduz também o valor do produto'' \parencite[p. 335-6]{Marx2014}.

\section{Cap. 14: ``O tempo de curso'' (ou tempo \textit{de circulação})}
Um fator importante do tempo de circulação do capital ``é a distância que separa o mercado do qual a mercadoria é vendida de seu local de produção'' \parencite[p. 344]{Marx2014}, pois, durante este período, ``o capital encontra-se imobilizado na condição de capital-mercadoria''. Dessa forma,
\begin{quote}
``O aperfeiçoamento dos meios de comunicação e de transporte abrevia o período de migração das mercadorias em termos absolutos, mas não suprime a diferença relativa, surgida dessa migração, entre os tempos de curso dos diferentes capitais-mercadoria, e tampouco entre diferentes frações do mesmo capital-mercadoria, que migram para diferentes mercados.'' \parencite[p. 344]{Marx2014}

``Ao mesmo tempo, com o desenvolvimento dos meios de transporte, não só se acelera a velocidade do deslocamento e, com isso, encurta-se temporalmente a distância espacial (...). Mas [também] quantidades sucessivas de mercadorias podem ser expedidas em períodos consecutivos mais curtos e, assim, chegar sucessivamente ao mercado [etc.]'' (p. 345)
\end{quote}

Dessa forma, torna-se possível que ``o tempo total de curso \[i.e. circulação\] \[seja\] abreviado e, assim, também a rotação'' (ibid.). Contudo, da mesma forma que abrevia os tempos de circulação de certas mercadorias — por exemplo, a expansão das ferrovias inglesas durante o século XIX impulsionou o comércio de carvão —, pode alongar os tempos de circulação de outras — como mercadorias que não podem, ou que não compensa, ser transportadas por estes meios de transporte mais velozes. 

Dessa forma, há uma relação reflexiva entre o desenvolvimento de centros urbanos e o desenvolvimento dos meios de transporte: centros produtivos induzem uma maior ``frequência de funcionamento dos meios de transporte'', com o que este centro produtivo consegue crescer mais, havendo um \textit{feedback loop}. Entrementes, tal desenvolvimento de centros produtivos ocorre \textit{em detrimento de demais localizações}, mesmo em detrimento de antigos polos produtivos — tal é a história do Ciclo colonial do ouro contra o ciclo do café, com que o interior de São Paulo ganhou maior prevalência econômica-política \textit{versus} o estado de Minas Gerais/Centro-Oeste.

``A rotação mais rápida do capital compensa aqui, em parte, o encarecimento de muitas condições da produção, dos terrenos onde se erguem as fábricas etc.'' \parencite[p. 346]{Marx2014}

\begin{quote}
``Se, por um lado, com o progresso da produção capitalista, o desenvolvimento dos meios de transporte e de comunicação abrevia o tempo de curso [i.e. circulação] para uma quantidade dada de mercadorias, esse mesmo progresso e a \textit{possibilidade} dada com o referido desenvolvimento provocam, inversamente, a \textit{necessidade} de trabalhar para mercados cada vez mais distantes — numa palavra, para o \textit{mercado mundial}. A massa das mercadorias que se encontram em deslocamento, sendo transportadas até pontos distantes, cresce enormemente e, com ela, tanto em termos absolutos quanto relativos, a parte do capital social que, constantemente e por períodos mais longos, encontra-se no estágio de capital-mercadoria, no interior do tempo de curso.'' (E, de fato, esta sociedade 'aparece como uma enorme coleção de mercadorias'!) ``Com isso também cresce, ao mesmo tempo, a parte da riqueza social que, em vez de servir de meio direto de produção, é investida em meios de transporte e de comunicação e no capital fixo e circulante requerido para o funcionamento desses meios.'' \parencite[p. 346, grifo meu]{Marx2014}
\end{quote}

Não só tais meios, como também os meios \textit{de pagamento}! O próprio Marx alude a esta problemática: supondo uma mercadoria cuja venda seja instantânea, i.e. ``que a mercadoria seja enviada logo após a realização do pedido e seja paga ao representante do produto no momento da entrega''; se o tempo de chegada da mercadoria é $t$ — no exemplo de Marx, $t=4 \text{ meses}$ —, então o tempo do refluxo dela como capital-dinheiro será $2t$ $=8\text{ meses}$. Portanto,
\begin{quote}
``o prolongamento do tempo provocado pela distância do mercado — tempo no qual o capital se desloca na forma de capital-mercadoria — provoca diretamente um retorno atrasado do dinheiro e, portanto, também atrasa a transformação do capital de [capital-dinheiro] em capital produtivo.'' \parencite[p. 349]{Marx2014}
\end{quote}

Hoje não há mais, em boa medida, esta problemática de se reaver o dinheiro de uma venda, não só por questão do crédito, como também pelo desenvolvimento dos próprios meios de pagamento, sistemas SWIFT, \textit{wire transfers}, Pix\footnote{\textbf{20260730}: Em 2026, Donald Trump tem \textit{upped the ante} contra o Pix do sistema brasileiro de pagamentos, com os motivos ``moralmente justificáveis'' de ``corrupção'', ``jogo sujo contra os \textit{players} financeiros'' e tantas baboseiras que não têm nada a ver com a questão. Deus nos ajude após a eleição de 2026... }, etc., além de instrumentos financeiros que garantem os preços de venda, como futuros etc. Mas de qualquer forma, tais problemáticas, e outras, induzem alguns capitalistas mais, e outros menos, a manter estoques produtivos.

\section{Cap. 15: ``Efeito do tempo de rotação sobre a grandeza do adiantamento de capital''}
\subsection{Primeira parte: análise de adiantamentos extras para continuidade de produção durante tempos de circulação}
Para as considerações abaixo, considera-se que o valor adiantado (médio) por semana é de £100. 

\subsubsection{Primeiro exemplo: rotação de 12 semanas (tempo de trabalho menor que tempo de circulação)}
\parencite[p. 356-7]{Marx2014} dá os seguintes períodos:
\begin{itemize}
\item Tempo de trabalho $t_{trab} = 9$ semanas
\item Período de não-trabalho $t_{\lnot trab}=0$, i.e. $t_{prod}=t_{trab}$
\item Tempo de circulação $t_{circ}=3$ semanas
\end{itemize}

Em vez de manter somente ``um'' capital operando por $9$ semanas e parado por $3$ semanas de circulação, um capitalista deseja manter ``um'' capital operando por $9$ semanas e, enquanto este leva $3$ semanas em circulação, manter \textit{outro capital} em período de trabalho ao mesmo tempo.

Para isso, Marx supõe que o adiantamento semanal é de $k=100$. Portanto, o adiantamento do período de trabalho é $9k$, e do período de circulação é $3k$. ``Originalmente'' — isto é, socialmente —, ele necessitaria somente $9k$ para fazer seu capital rodar, mas agora ele deseja investir para manter todas as semanas do ano operantes, portanto tendo de adiantar a mais. Com isso, ele obtém \textit{mais} do que uma rotação a cada 12 semanas.

Uma depicção disso é o seguinte:
!700

Cada quadrado representa três semanas; os quadrados hachurados em verde e em amarelo se tratam de trimestres \textit{ativamente em produção}: em verde são os que são empregados no primeiro período de trabalho, e o em amarelo é o que capital ``adicional'' que é empregado a partir do segundo período de trabalho; os quadrados hachurados em preto são períodos em que o capital produzido está \textit{em circulação}. Fica visível que, enquanto o \[que Marx chama de\] capital I está em circulação, o capital II começa sua rotação, e consegue acabar seu período de trabalho justamente porque o capital I completou sua rotação, trazendo de volta este valor adiantado, \textit{und so weiter}.

Vide descrição textual de Marx:
\begin{enumerate}
\item Primeira rotação: ``Primeiro período de trabalho de nove semanas: a rotação do capital aqui adiantado completa-se no início da 13ª semana'' (vide quarto quadrado verde na parte de baixo, começo da segunda rotação). ``Durante as últimas três semanas, funciona o capital adicional de £300'' (em amarelo), ``inaugurando o segundo período de trabalho de nove semanas.'' 
\item Segunda rotação: ``No começo da 13ª semana'' (após primeiro quadrado hachurado em preto), ``£900 retornaram e estão aptas a começar uma nova rotação. Mas o segundo período de trabalho já foi iniciado pelas £300 adicionais na 10ª semana'' (vide primeiro quadrado hachurado em amarelo). ``Das £900 iniciais foi liberada a soma de £300 para desempenhar o mesmo papel que o capital de £300 adicionado desempenhara no primeiro período de trabalho'' — i.e. vai ser agora um ``quadrado verde'' que começará um novo período de trabalho, após o fim do período de trabalho iniciado pelo ``quadrado amarelo''. E por aí vai.
\item Terceira rotação: similar. 
\item Quarta rotação: ``assim, o quarto período de rotação e o 5º \textit{de trabalho} [grifo meu] começam simultaneamente, com a 37ª semana do ano.''
\end{enumerate}

Note-se, portanto, que, enquanto o capital médio, adiantando $9k$, completou $3$ rotações — rodando, portanto, $27k$ —, este capital individual, adiantando $12k$, completou $3$ rotações e ainda está à espera da circulação ao final de mais uma rotação — rodou, portanto, $36k$, $9k$ dos quais ainda precisam passar pela circulação para serem vendidos.

\subsubsection{Segundo exemplo: rotação de 10 semanas (tempo de trabalho igual ao tempo de circulação)}
\parencite[p. 357]{Marx2014} :
\begin{itemize}
\item Tempo de trabalho $t_{trab}=5$ semanas
\item Período de não-trabalho $t_{\lnot trab}=0$, i.e. $t_{prod}=t_{trab}$
\item Tempo de circulação $t_{circ}=5$ semanas
\end{itemize}

Por simplicidade, pensemos que $5$ semanas sejam compostas de um só ``período'', a fim de que possamos pensar que uma rotação é composta de 1 parte de trabalho por 1 parte de circulação. Dessa forma, podemos enxergar as rotações alternadas da seguinte forma, em que cada quadrado represente $5$ semanas:
!400

Temos que, a cada duas rotações do capital médio, este capital individual terá feito duas rotações mais um período de trabalho completo (i.e. mais meia rotação). 

\subsubsection{Terceiro e quarto exemplos: rotação de 9 semanas (tempo de trabalho menor que tempo de circulação)}
\parencite[p. 371-3]{Marx2014}
\begin{itemize}
\item Tempo de trabalho $t_{trab}=3$ semanas
\item Período de não-trabalho $t_{\lnot trab}=0$, i.e. $t_{prod}=t_{trab}$
\item Tempo de circulação $t_{circ}=6$ semanas
\end{itemize}

Ou seja, adianta-se £900 para as 9 semanas, e os períodos de trabalho demandarão £300 cada.

Uma visualização deste cenário, em que cada quadrado representa $3$ semanas:
!Drawing 2026-07-31 12.46.06.svg

Note-se que a cada rotação (do capital médio), temos três movimentos simultâneos: o capital ``I'' faz uma rotação; o capital ``II'' faz $2/3$ de uma rotação (completa um período de trabalho e metade do tempo de circulação); e o capital ``III'' faz $1/3$ de uma rotação (somente completa o período de trabalho). 

Aqui é visível que ``Não ocorre um entrecruzamento ou entrelaçamento dos capitais'' \parencite[p. 373]{Marx2014}, devido à coincidência de que os tempos de produção e de circulação possuem um múltiplo em comum (no caso, $t_{circ}=2t_{trab}$). Faz-se necessário, portanto, analisar o caso bem mais comum de que estes tempos não sejam comensuráveis dessa forma.

\parencite[p. 373-6]{Marx2014}
\begin{itemize}
\item Tempo de trabalho $t_{trab}=4$ semanas
\item Período de não-trabalho $t_{\lnot trab}=0$, i.e. $t_{prod}=t_{trab}$
\item Tempo de circulação $t_{circ}=5$ semanas
\end{itemize}

Uma visualização desse caso, em que cada quadrado representa $1$ semana:
!Drawing 2026-07-31 13.08.19.svg

Note-se que o capital adicional (``capital III'') é empregado \textit{na última semana do tempo de circulação} do capital I, ou seja, permite começar o período de trabalho seguinte \textit{antecipadamente}, de forma que não haja semana com produção interrompida. 

\subsubsection{Referências}
